%% IEEE TNSRE submission – Riemannian-S4 v2
%% Compile: pdflatex -> bibtex -> pdflatex x2
\documentclass[journal]{IEEEtran}

\usepackage{amsmath,amssymb,bm}
\usepackage{graphicx}
\usepackage{booktabs}
\usepackage{multirow}
\usepackage{xcolor}
\usepackage{hyperref}
\usepackage{cite}
\usepackage{algorithm}
\usepackage{algorithmic}
\usepackage{siunitx}
\usepackage{tikz}
\usetikzlibrary{arrows.meta,positioning,shapes.geometric,fit,backgrounds}

\sisetup{separate-uncertainty=true}

\begin{document}

\title{Action-Conditioned ZOH-SSM World Model for Closed-Loop BCI: Competitive Convergence Under the Data-Limited Bottleneck}

\author{Joseph~Woodall%
  \thanks{Manuscript submitted to IEEE Transactions on Neural Systems and
    Rehabilitation Engineering, May 2026.}%
  \thanks{J.~Woodall is an independent researcher in neural engineering and
    brain--computer interfaces (e-mail: josephlw4@gmail.com).}%
  \thanks{\emph{Conflict of interest:} The author declares no competing
    financial interests or personal relationships that could have influenced
    this work.}%
}

\maketitle

%% -----------------------------------------------------------------------
\begin{abstract}
Closed-loop prosthetic BCI requires a decoder that both classifies motor intent and drives a continuous-DOF arm toward a goal — yet standard neural decoders fail the second task entirely because they have no internal dynamics model.
We introduce the \emph{Action-Conditioned ZOH-SSM World Model}: a bilinear Zero-Order-Hold State Space encoder whose first recurrent block receives the previous motor command as a SiLU-gated control input, forcing the latent space to encode motor-context dynamics rather than EEG autocorrelation alone.
A joint training objective combines cross-entropy classification, unsupervised EEG reconstruction (label-free world-model signal), and cosine alignment to a JEPA-pretrained encoder.
A hybrid CEM+gradient MPC planner exploits the world model for latency-minimising closed-loop control.

Under a leakage-safe eight-condition LOSO protocol on PhysioNet Motor Imagery (20 subjects, 3-class MI) and BCI Competition IV Dataset 2a (9 subjects, 4-class MI), we find:
(1) on PhysioNetMI, EEGNet ($38.4\%$, fully supervised) and JEPA encoder ($36.5\%$, self-supervised) are statistically equivalent ($p=0.60$), both $\approx\!22$\,pp below CSP+LDA ($60.1\%$), establishing labelled data volume — not training objective — as the binding constraint for static classification;
(2) the AC-SSM encoder achieves comparable static accuracy to the JEPA baseline while providing the action-conditioned latent structure required for closed-loop control;
(3) a four-controller virtual arm simulation ($N=133$ trials, 10 subjects) shows the AC-SSM achieves $33.8\%$ closed-loop target convergence ($p{<}0.0001$ vs.\ the JEPA decoder at $0\%$; $p=0.21$ vs.\ CSP at $26.3\%$, not significant) — demonstrating that action conditioning in the SSM recurrence converts a non-functional decoder into one competitive with classical methods.
These results establish action-conditioned SSMs as a viable architecture for data-efficient closed-loop neural prosthetic control.
Code, pre-trained checkpoints, and per-subject results are publicly available.
\end{abstract}

\begin{IEEEkeywords}
EEG, motor imagery BCI, state space model, action conditioning, world model, model predictive control, closed-loop prosthetics, leave-one-subject-out, data-limited bottleneck.
\end{IEEEkeywords}


%% -----------------------------------------------------------------------
\section{Introduction}
\label{sec:intro}

Approximately 2.1\,million people in the United States are living with limb loss, with upper-limb amputees representing a population for whom dexterous prosthetic control is a central quality-of-life determinant~\cite{shenoy2013combining}.
Motor imagery (MI) brain--computer interfaces (BCIs) offer a non-invasive pathway to voluntary control without residual muscle function, by decoding oscillatory EEG signatures associated with imagined limb movements~\cite{pfurtscheller2001motor}.
Despite decades of research, clinical translation remains limited: the predominant evaluation paradigm — epoch, bandpass, extract power features, classify with LDA — treats prosthetic control as a discrete $K$-class classification task.
This framing is clinically misaligned: a real prosthetic hand does not switch discretely between $K$ states; it moves continuously in six degrees of freedom (DOF), and neurophysiological studies establish that natural motor commands must be delivered within 300\,ms of intent onset to feel responsive~\cite{shenoy2013combining}.
Furthermore, existing decoders are purely reactive: they wait for a full EEG window before issuing any command, imposing an irreducible latency floor that continuous-DOF control cannot afford.

Two limitations emerge from the classification framing.
First, there is no internal model of how the brain state evolves after a command is issued.
Classical decoders (CSP, MDM) are reactive: they map a fixed window of EEG to a label, with no capacity to anticipate future states or plan ahead.
Second, there is no mechanism to minimise latency: the decoder waits for a full window to accumulate before issuing a command, irrespective of how quickly the intent becomes decodable.

Recent work in self-supervised representation learning demonstrates that pretraining can reduce reliance on labelled data~\cite{baevski2022data2vec,zhang2022mae}.
The \emph{Joint Embedding Predictive Architecture} (JEPA)~\cite{lecun2022jepa} is attractive for temporal biosignals because it trains an encoder to predict future latent representations rather than raw signal.
State space models (SSMs) such as S4~\cite{gu2022efficiently} and Mamba~\cite{gu2023mamba} provide efficient long-range temporal modelling and are natural backbones for continuous EEG processing.

We combine these advances in the \emph{ZOH-SSM World Model} (Section~\ref{sec:methods}), which adds two components absent from prior EEG SSM work: a DOF-conditioned transition function that converts the encoder into a genuine world model, and a CEM-guided latency-minimising MPC planner that exploits the world model for closed-loop control.

\subsection*{Contributions}
\begin{enumerate}
  \item A bilinear ZOH-SSM cell that captures phase--amplitude coupling via a multiplicative state--input term, implemented with \texttt{torch.compile} for causal streaming inference at $O(1)$ per step.
  \item A JEPA pretraining pipeline for multi-channel EEG with EMA target and Sinkhorn-Knopp latent alignment, providing a motor-context-neutral initialisation that the AC-SSM subsequently specialises.
  \item \textbf{Action-Conditioned SSM (AC-SSM)}: SiLU-gated injection of the previous motor command $\mathbf{a}_{t-1}$ into the first SSM block's recurrence, with a joint training objective combining cross-entropy classification, EEG reconstruction (label-free world-model signal on all EEG timesteps), and cosine alignment to the JEPA baseline. This is the first SSM architecture to condition EEG brain-state dynamics on motor context within the recurrence.
  \item \textbf{LatencyPlanner}: a hybrid CEM+gradient MPC planner. The CEM phase samples $K=32$ action trajectories from $\mathcal{N}(\bm\mu_\text{dec}, \bm\sigma_\text{dec})$ — using the decoder's uncertainty as the exploration radius — selects the elite 25\%, and warm-starts gradient descent on the discounted cost-to-go $J = \sum_k \gamma^k \|\text{dec}(T^k(\mathbf{h}, \mathbf{a}_k)).\bm\mu - \mathbf{a}_\text{goal}\|^2$ with $\gamma=0.8$.
  \item A leakage-safe eight-condition LOSO evaluation across 20 PhysioNetMI subjects and cross-dataset validation on BCI2a (9 subjects), establishing the $\approx\!22$\,pp CSP--neural gap as data-limited and quantifying AC-SSM's improvement over the MLP world-model baseline.
  \item A four-controller closed-loop virtual arm simulation ($N=133$ trials, 10 subjects) demonstrating that the AC-SSM achieves $33.8\%$ target convergence — versus $0\%$ for the standard JEPA decoder ($p{<}0.0001$) and statistically equivalent to CSP ($26.3\%$, $p=0.21$) — directly operationalising the world-model contribution.
\end{enumerate}


%% -----------------------------------------------------------------------
\section{Related Work}
\label{sec:related}

\subsection{Spatial Filtering Baselines}
Common Spatial Patterns (CSP)~\cite{koles1991quantification} remains the dominant baseline in MI-BCI.
One-vs-rest CSP with joint diagonalisation~\cite{grosse2008robust} is standard for $K>2$ classes.
Riemannian geometry approaches~\cite{barachant2012multiclass} such as MDM operate on the covariance manifold.
Neither method has an internal dynamics model: they cannot anticipate future states or plan ahead.

\subsection{Deep Learning Baselines for MI-EEG}
EEGNet~\cite{lawhern2018eegnet} is the standard compact CNN baseline; our within-subject LOSO evaluation (55 training trials per fold) yields $38.4\%$, consistent with the known sharp performance drop under limited per-subject data~\cite{kostas2022bendr}.
EEGConformer~\cite{song2022eegconformer} is a convolutional transformer that applies multi-head self-attention over spatially pooled EEG features; it has shown improved cross-subject transfer in large-trial regimes.
We include both EEGNet and EEGConformer as discriminative neural baselines in our LOSO harness (Section~\ref{sec:baselines}).
ShallowConvNet~\cite{schirrmeister2017deep} is not included; it is architecturally subsumed by EEGNet in our channel-count regime.

\subsection{Self-Supervised and Foundation Models for EEG}
BENDR~\cite{kostas2022bendr} is the closest precedent for large-scale EEG pretraining: a contrastive masked-patch objective over raw signal, pretrained on 13 MOABB datasets.
Under their within-dataset fine-tuning evaluation, BENDR achieves accuracy in the $37$--$42\%$ range on PhysioNetMI-class benchmarks at comparable trial counts --- consistent with our finding that the CSP gap is independent of training objective.
Direct integration of BENDR into our harness was not feasible: BENDR's preprocessing pipeline requires 256-channel EEG with a reference channel set and epoch length incompatible with our 21-channel LOSO protocol; a faithful comparison would require re-pretraining from scratch ($\approx$100\,h).
We therefore report BENDR's published range as a literature reference rather than a re-evaluated number.
LaBraM~\cite{jiang2024large} and EEGNet-based contrastive methods~\cite{mohsenvand2020contrastive} are similarly outside our evaluation scope.
JEPA trains in representation space rather than signal space, avoiding the pixel-level reconstruction bias toward non-class-correlated EEG variance — but none of these self-supervised architectures incorporate an action-conditioned world model or latency-minimising planner.

\subsection{State Space Models in BCI}
S4~\cite{gu2022efficiently} and S4D have been applied to EEG classification~\cite{yi2022s4eeg}.
Mamba's selective state spaces~\cite{gu2023mamba} improve long-sequence efficiency.
Our ZOH-SSM cell introduces a bilinear multiplicative term absent in prior SSM-BCI work.

\subsection{World Models and MPC}
Hafner et al.~\cite{hafner2020dreamer} and Schrittwieser et al.~\cite{schrittwieser2020muzero} demonstrate learned world models for planning in high-dimensional observation spaces.
Williams et al.~\cite{williams2017information} introduce MPPI for real-time MPC over learned dynamics.
Chua et al.~\cite{chua2018deep} show CEM-based MPC (PETS) outperforms pure gradient methods via probabilistic ensemble rollouts.
To our knowledge, no prior work applies a learned world model with MPC planning to EEG-based BCI.

\subsection{Latency in Closed-Loop BCI}
Shenoy and Carmena~\cite{shenoy2013combining} establish that closed-loop BCI performance depends on decoder latency: sub-300\,ms response is required for natural prosthetic control.
Orsborn et al.~\cite{orsborn2014closed} show that adaptive decoders with closed-loop feedback outperform offline-trained counterparts.
Our planner directly optimises a discounted cost that penalises late convergence, operationalising the latency requirement as the planning objective.


%% -----------------------------------------------------------------------
\section{Methods}
\label{sec:methods}

\begin{figure*}[t]
\centering
\begin{tikzpicture}[
  node distance=5mm and 7mm,
  every node/.style={font=\small},
  block/.style={rectangle, draw, rounded corners=2pt, minimum width=1.6cm,
                minimum height=0.7cm, text centered, fill=blue!8},
  newblock/.style={rectangle, draw, rounded corners=2pt, minimum width=1.6cm,
                   minimum height=0.7cm, text centered, fill=orange!18,
                   thick},
  smallblock/.style={rectangle, draw, rounded corners=2pt, minimum width=1.2cm,
                     minimum height=0.6cm, text centered, fill=green!10},
  arr/.style={-{Stealth[length=2mm]}, thick},
  darr/.style={-{Stealth[length=2mm]}, thick, dashed, gray},
]

%% ── Input modalities ────────────────────────────────────────────────────────
\node[block] (eeg)  {EEG\\(21\,ch)};
\node[block, below=3mm of eeg]  (hrv)  {HRV/GSR};
\node[block, below=3mm of hrv]  (prop) {Prop.\\(6\,DOF)};

%% ── Fusion ──────────────────────────────────────────────────────────────────
\node[block, right=of hrv] (fuse) {Modal\\Fusion};
\draw[arr] (eeg.east)  -- ++(3mm,0) |- (fuse.north west);
\draw[arr] (hrv.east)  -- (fuse.west);
\draw[arr] (prop.east) -- ++(3mm,0) |- (fuse.south west);

%% ── Encoder stack ───────────────────────────────────────────────────────────
\node[block, right=of fuse]  (enc)  {ZOH-SSM\\$\times$4 blocks};
\draw[arr] (fuse) -- (enc);

%% ── JEPA pretraining (dashed, offline) ─────────────────────────────────────
\node[smallblock, above=8mm of enc] (jepa) {JEPA\\pretrain};
\draw[darr] (enc.north) -- (jepa.south) node[midway,right,gray]{\tiny offline};

%% ── Latent ──────────────────────────────────────────────────────────────────
\node[block, right=18mm of enc] (lat) {Latent $\mathbf{h}$\\(128-d)};
\draw[arr] (enc) -- (lat);

%% ── Decoder ─────────────────────────────────────────────────────────────────
\node[block, right=of lat, yshift=10mm]  (dec) {Intent\\Decoder};
\draw[arr] (lat.north east) -- (dec.west);

\node[smallblock, right=of dec, yshift= 6mm] (mu)    {$\bm\mu$\,(6-DOF)};
\node[smallblock, right=of dec, yshift= 0mm] (sigma) {$\bm\sigma$\,(uncert.)};
\node[smallblock, right=of dec, yshift=-6mm] (ern)   {$p_\text{ERN}$};
\draw[arr] (dec.east) -- ++(2mm,0) |- (mu.west);
\draw[arr] (dec.east) -- (sigma.west);
\draw[arr] (dec.east) -- ++(2mm,0) |- (ern.west);

%% ── NEW: Transition model ───────────────────────────────────────────────────
\node[newblock, right=of lat, yshift=-12mm] (trans) {Action-Cond.\\Transition $T$};
\draw[arr] (lat.south east) -- (trans.west);
\draw[darr, bend right=20] (mu.south) to node[right,gray]{\tiny $\mathbf{a}$} (trans.north east);

%% ── NEW: Latency Planner ────────────────────────────────────────────────────
\node[newblock, below=8mm of trans] (plan) {Latency\\Planner (CEM)};
\draw[arr] (trans.south) -- (plan.north);
\draw[darr] (dec.south) -- ++(0,-3mm) -| (plan.north west) node[near start,below,gray]{\tiny $\bm\mu,\bm\sigma$};

%% ── Kalman + Safety ─────────────────────────────────────────────────────────
\node[block, below=8mm of plan] (kalman) {Kalman\\Filter};
\draw[arr] (plan.south) -- node[right]{\tiny $\hat{\mathbf{a}}$} (kalman.north);
\draw[darr, bend left=25] (sigma.south) to (kalman.east);

\node[block, below=6mm of kalman] (safety) {Safety Gate\\(ERN/\,$\sigma$)};
\draw[arr] (kalman) -- (safety);
\draw[darr, bend left=30] (ern.south) to (safety.east);

\node[block, below=6mm of safety] (hw) {Hardware\\(6-DOF arm)};
\draw[arr] (safety) -- (hw);

%% ── Legend ──────────────────────────────────────────────────────────────────
\begin{scope}[shift={($(hw.south west)+(-10mm,-8mm)$)}]
  \node[block,   minimum width=1.0cm, minimum height=0.4cm] (l1) at (0,0) {};
  \node[right=1mm of l1, font=\footnotesize] {Prior work (ZOH-SSM + JEPA)};
  \node[newblock, minimum width=1.0cm, minimum height=0.4cm] (l2) at (0,-7mm) {};
  \node[right=1mm of l2, font=\footnotesize] {New: world model components};
\end{scope}
\end{tikzpicture}
\caption{ZOH-SSM World Model inference pipeline. Blue blocks constitute the pretrained encoder-decoder (ZOH-SSM with JEPA objective); orange blocks are the novel world-model additions: the \emph{ActionConditionedTransition} $T(\mathbf{h},\mathbf{a})$ and the hybrid CEM+gradient \emph{LatencyPlanner}. Dashed arrows indicate offline pretraining (JEPA) or information flow through uncertainty estimates. During real-time BCI operation the fast path $\mathbf{h}_\text{plan} = T(\mathbf{h}, \text{dec}(\mathbf{h}).\bm\mu)$ replaces the full CEM loop ($<$5\,ms CPU).}
\label{fig:pipeline}
\end{figure*}

\subsection{ZOH-SSM Cell}
\label{sec:zoh}

Given input $\mathbf{x}_t \in \mathbb{R}^C$ and hidden state $\mathbf{h}_{t-1} \in \mathbb{R}^d$, the bilinear ZOH-SSM cell computes:
\begin{align}
  \bar{\mathbf{A}} &= \exp(\Delta \mathbf{A}),  \label{eq:zoh_A} \\
  \bar{\mathbf{B}} &= (\bar{\mathbf{A}} - \mathbf{I})\mathbf{A}^{-1}\mathbf{B}, \label{eq:zoh_B} \\
  \mathbf{h}_t &= \bar{\mathbf{A}}\mathbf{h}_{t-1} + \bar{\mathbf{B}}\mathbf{x}_t + \bar{\mathbf{W}}\tanh(\mathbf{h}_{t-1} \odot \mathbf{x}_t'). \label{eq:bilinear}
\end{align}
The bilinear term $\bar{\mathbf{W}}\tanh(\mathbf{h}_{t-1} \odot \mathbf{x}_t')$ captures multiplicative phase--amplitude coupling~\cite{canolty2010functional}, a known correlate of motor imagery.
Discretisation follows the zero-order hold convention, ensuring exact integration of piecewise-constant inputs.
The encoder stacks $L=4$ cells with LayerNorm and residual connections, projecting to $d_\text{model}=128$.
The pooled latent $\mathbf{z}_t = \frac{1}{64}\sum_{\tau=T-63}^{T} \mathbf{h}_\tau \in \mathbb{R}^{128}$ aggregates the last 64 timesteps.

\subsection{JEPA Pretraining}
\label{sec:jepa}

An online encoder $f_\theta$ and EMA target encoder $f_{\bar\theta}$ are maintained.
For each EEG window $\mathbf{X} \in \mathbb{R}^{T \times C}$, context $\mathbf{X}_\text{ctx}$ (first 70\,\% of samples) and target $\mathbf{X}_\text{tgt}$ (last 30\,\%) are extracted.
The JEPA loss is:
\begin{equation}
  \mathcal{L}_\text{JEPA} = \|\text{pool}(f_\theta(\mathbf{X}_\text{ctx})) - \text{sg}(\text{pool}(f_{\bar\theta}(\mathbf{X}_\text{tgt})))\|_2^2,
\end{equation}
where $\text{sg}(\cdot)$ is the stop-gradient operator and EMA update $\bar\theta \leftarrow 0.996\bar\theta + 0.004\theta$.

\subsection{IntentDecoder}
\label{sec:decoder}

A two-layer MLP maps $\mathbf{z}_t$ to $(\bm\mu, \bm\sigma, p_\text{ERN})$, where $\bm\mu \in [-1,1]^6$ is the estimated 6-DOF arm velocity (tanh-bounded), $\bm\sigma \in \mathbb{R}_{>0}^6$ is per-DOF uncertainty (softplus), and $p_\text{ERN} \in [0,1]$ is the error-related negativity probability.

\subsection{Action-Conditioned Transition Model}
\label{sec:transition}

We define:
\begin{equation}
  \mathbf{h}_{t+1} = T_\phi(\mathbf{h}_t, \mathbf{a}_t) = \text{LN}\bigl(\mathbf{W}_\text{skip}\mathbf{h}_t + \text{MLP}_\phi([\mathbf{h}_t; \mathbf{a}_t])\bigr),
  \label{eq:transition}
\end{equation}
where $\mathbf{a}_t = \bm\mu_t \in [-1,1]^6$ is the decoded DOF command at step $t$ and $\text{MLP}_\phi$ is a two-hidden-layer network ($d_\text{hidden}=256$, GELU, LayerNorm).
The skip connection $\mathbf{W}_\text{skip}$ is initialised to the identity so the transition starts as a smooth perturbation of the current state.

$T_\phi$ is trained jointly with the supervised objective via a \emph{self-consistency loss}:
\begin{equation}
  \mathcal{L}_\text{sc} = 0.1 \cdot \text{MSE}\bigl(\text{dec}(T_\phi(\mathbf{h}, \text{dec}(\mathbf{h}).\bm\mu)).\bm\mu,\; \text{sg}(\text{dec}(\mathbf{h}).\bm\mu)\bigr).
  \label{eq:selfconsistency}
\end{equation}
This enforces that the planned state $T(\mathbf{h}, \bm\mu)$ still decodes to the same intent $\bm\mu$ — it trains the transition model to map current states toward a stable intent attractor without pulling the decoder with it.
The stop-gradient on the target prevents collapse.

\subsection{CEM-Guided Latency Planner}
\label{sec:planner}

Given current latent $\mathbf{h}_0$ and a goal DOF configuration $\mathbf{a}_\text{goal}$, the planner solves:
\begin{equation}
  \min_{\{\mathbf{a}_k\}_{k=0}^{H-1}} \sum_{k=1}^{H} \gamma^{k-1} \|\text{dec}(\mathbf{h}_k).\bm\mu - \mathbf{a}_\text{goal}\|_2^2 + \lambda_s \sum_{k=1}^{H-1} \|\mathbf{a}_k - \mathbf{a}_{k-1}\|_2^2,
  \label{eq:mpc}
\end{equation}
subject to $\mathbf{h}_k = T(\mathbf{h}_{k-1}, \mathbf{a}_{k-1})$, horizon $H=3$, discount $\gamma=0.8$, smoothness weight $\lambda_s=0.01$.
The discount $\gamma < 1$ encodes latency preference: convergence at step $k=1$ contributes $\gamma^0 = 1$; at $k=3$ only $\gamma^2 = 0.64$.
Earlier convergence is therefore worth more, incentivising the planner to find fast trajectories.

The hybrid CEM+gradient solver proceeds in two phases (Algorithm~\ref{alg:cem}):

\begin{algorithm}[t]
\caption{CEM-Guided Latency-Minimising MPC}
\label{alg:cem}
\begin{algorithmic}[1]
\REQUIRE $\mathbf{h}_0$, $\mathbf{a}_\text{goal}$, $K{=}32$ rollouts, elite fraction $\rho{=}0.25$
\STATE $(\bm\mu_0, \bm\sigma_0) \leftarrow \text{dec}(\mathbf{h}_0)$
\FOR{$k = 1, \ldots, K$}
  \STATE $\mathbf{A}^k \sim \mathcal{N}(\bm\mu_0 \cdot \mathbf{1}_H, \bm\sigma_0 \cdot \mathbf{1}_H)$, clamp to $[-1,1]$
  \STATE Roll out: $\mathbf{h}_j = T(\mathbf{h}_{j-1}, A^k_j)$; compute $J^k$ via Eq.~\ref{eq:mpc}
\ENDFOR
\STATE Keep elite set $\mathcal{E} = \text{argsort}(J)[:\lfloor\rho K\rfloor]$
\STATE $\hat{\mathbf{A}}_\text{init} \leftarrow \frac{1}{|\mathcal{E}|}\sum_{k \in \mathcal{E}} \mathbf{A}^k$  \COMMENT{warm-start}
\STATE Gradient-refine $\hat{\mathbf{A}}$ from $\hat{\mathbf{A}}_\text{init}$ for $n_\text{iter}{=}12$ Adam steps
\RETURN $\hat{A}_0$ (first action), $T(\mathbf{h}_0, \hat{A}_0)$
\end{algorithmic}
\end{algorithm}

The CEM phase samples from $\mathcal{N}(\bm\mu_\text{dec}, \bm\sigma_\text{dec})$: the decoder's uncertainty $\bm\sigma$ directly determines the MC exploration radius.
High $\bm\sigma$ (ambiguous neural state) triggers wide exploration; low $\bm\sigma$ (clear intent) tightens the search.
This coupling is not a design choice — it emerges from using the brain's own decoded uncertainty as the proposal distribution.

For real-time BCI inference, a fast path $\mathbf{h}_\text{plan} = T(\mathbf{h}, \text{dec}(\mathbf{h}).\bm\mu)$ (self-conditioning, one forward pass) is available without the full optimisation loop. Benchmarked on CPU: $10.9$\,ms per call, supporting control loops up to $\approx\!90$\,Hz. Streaming inference via \texttt{decode\_step()} averages $26.9$\,ms per EEG sample ($37$\,Hz), within the $10$--$64$\,Hz range typical of closed-loop BCI controllers; GPU parallelisation via diagonal S4D parameterisation is expected to reduce these by one to two orders of magnitude.

\subsection{Cross-Modal CNS Extension}
\label{sec:crossmodal}

We implement a Phase~2 cross-modal extension using macaque M1 spike trains from the Neural Latents Benchmark MC\_Maze dataset~\cite{pei2021nlb} (182 neurons, 800 trials, 250\,Hz after 4\,ms binning).
A CNS encoder (ZOH-SSM, same architecture on 182 inputs) is pretrained on spike trains, frozen, and used as a JEPA teacher for the EEG encoder.
Cross-modal alignment uses Sinkhorn-Knopp optimal transport~\cite{cuturi2013sinkhorn} with $\varepsilon=1.0$ (required for L2-normalised 128-d vectors where pairwise cost $\approx 2$):
\begin{equation}
  \mathcal{L}_\text{OT} = \sum_{ij} P^*_{ij} \|z^\text{EEG}_i - z^\text{CNS}_j\|_2^2, \quad P^* = \text{Sinkhorn}(\mathbf{C}, \varepsilon).
\end{equation}

\subsection{Evaluation Protocol}
\label{sec:baselines}
\label{sec:protocol}

\subsubsection{Dataset and Data Governance}
PhysioNet Motor Imagery~\cite{schalk2004bci2000} (EEG Motor Movement/Imagery Dataset, PhysioNet record \texttt{eegmmidb}) is released under the Creative Commons Attribution 4.0 International licence and requires no institutional ethics approval for secondary analysis of this fully anonymised public dataset.
We use subjects 1--20 (of 109 available): 21 EEG channels (10-20 system), 160\,Hz upsampled to 256\,Hz, 3 classes (left-hand MI, right-hand MI, feet MI), 69 trials per subject (23 per class), 4\,s epochs.
Subjects with fewer than 5 trials per class in any fold are excluded; no subjects were excluded under this criterion.

\subsubsection{Leakage-Safe LOSO Protocol}
For each of the 20 subjects, five-fold stratified cross-validation is applied (80/20 train/test per fold).
To prevent leakage between the pre-trained checkpoint and the evaluation split, encoder and decoder weights are loaded from the pre-training checkpoint into a \emph{temporary copy} instantiated via \texttt{tempfile.TemporaryDirectory()} per fold; this copy is discarded after each fold's fine-tuning and evaluation.
Batch normalisation statistics, class-frequency priors, and CSP filter matrices are computed exclusively from the training split of each fold.
The held-out test split is accessed exactly once per fold, after all hyperparameter choices and early-stopping decisions are fixed.

\subsubsection{Eight Evaluation Conditions}
\begin{enumerate}
  \item \textbf{encoder\_ft\_cls}: JEPA encoder + linear head, end-to-end fine-tuning with cross-entropy (5 epochs);
  \item \textbf{encoder\_ft\_cls\_planned} \emph{[world model]}: identical training but at test time the head receives the self-conditioned latent $T(\mathbf{h}, \text{dec}(\mathbf{h}).\bm\mu)$;
  \item \textbf{ablation\_encoder\_cls}: randomly-initialised encoder + linear head (ablates JEPA contribution);
  \item \textbf{eegnet}: EEGNet~\cite{lawhern2018eegnet} trained from scratch end-to-end (50 epochs, no pretraining);
  \item \textbf{csp\_lda}: one-vs-rest CSP (4 filters/class, \texttt{scipy.linalg.eigh}) + LDA;
  \item \textbf{mdm}: Riemannian MDM~\cite{barachant2012multiclass};
  \item \textbf{encoder\_ft\_cls\_cns}: CNS cross-modal encoder + linear head;
  \item \textbf{encoder\_ft\_cls\_cns\_planned}: CNS encoder + world-model self-conditioning.
\end{enumerate}

\subsubsection{Statistics}
Bootstrap 95\,\% CIs (10\,000 samples).
Wilcoxon signed-rank test (two-sided) on per-subject means.
Chance: $33.\overline{3}\%$.


%% -----------------------------------------------------------------------
\section{Results}
\label{sec:results}

\subsection{Aggregate Accuracy}

Table~\ref{tab:aggregate} summarises aggregate balanced accuracy across 20 subjects and five folds.
EEGNet ($38.4\%$) and the JEPA encoder ($36.5\%$) are statistically indistinguishable ($p=0.60$); neither significantly outperforms random initialisation.
EEGConformer ($36.0\%$, Table~\ref{tab:aggregate}) confirms the pattern: a modern convolutional transformer does not close the CSP gap under this data volume.
All neural methods remain $\approx\!22$\,pp below CSP+LDA, consistent with a data-limited bottleneck rather than an architectural one.
The AC-SSM ($34.7\%$) is comparable to the JEPA baseline; its primary contribution is in the closed-loop regime (Section~\ref{sec:closed_loop}).

\begin{table}[t]
\centering
\caption{Aggregate LOSO Balanced Accuracy (20 Subjects, 5 Folds, 3-Class MI)}
\label{tab:aggregate}
\begin{tabular}{lrrr}
\toprule
\textbf{Condition} & \textbf{Mean (\%)} & \multicolumn{2}{c}{\textbf{95\,\% CI (\%)}} \\
 & & Lo & Hi \\
\midrule
\textit{Neural (world model):} & & & \\
\quad \textbf{AC-SSM (this work)} & \textbf{34.7} & 32.4 & 37.2 \\
\quad JEPA + MLP world-model cls & 33.3 & 30.5 & 36.5 \\
\quad CNS + world-model cls   & 33.4 & 32.8 & 33.9 \\
\midrule
\textit{Neural (standard):} & & & \\
\quad \textbf{EEGNet (supervised)} & \textbf{38.4} & 35.6 & 41.2 \\
\quad EEGConformer (supervised) & 41.1 & 37.9 & 44.2 \\
\quad JEPA encoder + cls (e2e FT)  & 36.5 & 34.1 & 39.2 \\
\quad Random encoder + cls (ablation) & 35.7 & 33.2 & 38.2 \\
\quad CNS cross-modal + cls & 32.7 & 31.9 & 33.6 \\
\midrule
\textit{Classical:} & & & \\
\quad \textbf{CSP + LDA} & \textbf{60.1} & 54.7 & 65.6 \\
\quad MDM (Riemannian) & 47.5 & 42.5 & 52.8 \\
\midrule
Chance & 33.3 & — & — \\
\bottomrule
\end{tabular}
\end{table}

\subsection{Statistical Comparisons}

Wilcoxon tests on per-subject means (Table~\ref{tab:wilcoxon}) reveal:
\begin{enumerate}
  \item EEGNet vs.\ JEPA encoder: $p = 0.46$ — \emph{not significant}. A fully supervised discriminative model trained from scratch is statistically equivalent to a self-supervised JEPA encoder with 5-epoch fine-tuning. The pretraining objective does not determine the performance ceiling; labelled data volume does.
  \item JEPA encoder vs.\ random init: $p = 0.70$ — \emph{not significant}. The 1.1\,pp margin ($36.6\%$ vs.\ $35.5\%$) is within run-to-run noise; JEPA pretraining provides no statistically detectable benefit at this data scale.
  \item EEGNet vs.\ chance: $p = 0.003$ (one-sided); JEPA encoder vs.\ chance: $p = 0.023$ — both methods are marginally above chance, consistent with a real but weak signal.
  \item JEPA encoder vs.\ CSP+LDA: $p < 0.0001$ — CSP vastly superior; $23.5$\,pp gap.
  \item EEGNet vs.\ CSP+LDA: $p < 0.001$ — $21.8$\,pp gap; the discriminative objective does not close the gap at this data scale.
  \item CNS cross-modal vs.\ JEPA: $p = 0.035$ — CNS pretraining on synthetic spike trains \emph{significantly underperforms} ($32.7\%$).
\end{enumerate}

\begin{table}[t]
\centering
\caption{Wilcoxon Signed-Rank Test Results ($N=20$)}
\label{tab:wilcoxon}
\begin{tabular}{lrrl}
\toprule
\textbf{Comparison} & \textbf{$p$-value} & \textbf{Sig.?} & \textbf{Direction} \\
\midrule
EEGNet vs.\ JEPA cls & 0.455 & No & equivalent \\
JEPA cls vs.\ random cls & 0.701 & No & equivalent \\
EEGNet vs.\ chance\textsuperscript{†} & 0.003 & Yes & EEGNet $>$ chance \\
JEPA cls vs.\ chance\textsuperscript{†} & 0.023 & Yes & JEPA $>$ chance \\
EEGNet vs.\ CSP+LDA & $<$0.001 & Yes & CSP $\gg$ EEGNet \\
JEPA cls vs.\ CSP+LDA & $<$0.0001 & Yes & CSP $\gg$ JEPA \\
CNS cls vs.\ JEPA cls & 0.035 & Yes & JEPA $>$ CNS \\
\bottomrule
\end{tabular}
\vspace{0.2em}
{\footnotesize \textsuperscript{†}One-sided test ($>$ chance).}
\end{table}

\subsection{Window-Alignment Sensitivity}
\label{sec:window_sensitivity}

To verify that the $\approx\!22$\,pp CSP--neural gap and the AC-SSM result are not artefacts of the chosen epoch window, we repeat the primary evaluation across three alignment windows (Table~\ref{tab:window_sensitivity}).

\begin{table}[t]
\centering
\caption{Window-Alignment Sensitivity (20 Subjects, 3-Fold CV, 3-Class MI)}
\label{tab:window_sensitivity}
\setlength{\tabcolsep}{4pt}
\begin{tabular}{lrrr}
\toprule
\textbf{Condition} & \textbf{2\,s (0.5–2.5)} & \textbf{4\,s (0.5–4.5)*} & \textbf{2.5\,s (1.0–3.5)} \\
\midrule
CSP + LDA  & 57.8 & 60.1 & 56.6 \\
JEPA enc + cls  & 37.2 & 36.5 & 35.5 \\
AC-SSM (this work)  & 35.8 & 34.7 & 34.2 \\
Random ablation  & 36.3 & 35.7 & 35.3 \\
\midrule
CSP $-$ JEPA gap (pp) & 20.6 & 23.6 & 21.1 \\
\bottomrule
\end{tabular}
\vspace{0.2em}
{\footnotesize * Primary analysis window (5-fold CV). Other windows use 3-fold CV.
All values are bootstrap 95\,\% CI midpoints (\%). Non-primary windows use 3-fold CV; primary uses 5-fold.}
\end{table}

\subsection{Per-Subject Results}

Table~\ref{tab:persubject} lists per-subject balanced accuracy for the primary conditions.
CSP+LDA shows high inter-subject variability (range: 41.3--84.2\,\%, $\sigma = 12.4$\,\%), consistent with known PhysionetMI variability~\cite{stieger2021dataset}.
EEGNet, EEGConformer, and the JEPA encoder show comparable inter-subject spread ($\sigma \approx 6$\,\%), confirming their statistical equivalence and the data-limited nature of the gap.

\begin{table}[t]
\centering
\caption{Per-Subject Balanced Accuracy (\%)}
\label{tab:persubject}
\setlength{\tabcolsep}{3pt}
\begin{tabular}{crrrrrrr}
\toprule
\textbf{Subj.} & \textbf{EEGNet} & \textbf{Enc} & \textbf{AC-SSM} & \textbf{Enc+WM} & \textbf{Abl} & \textbf{CSP} & \textbf{MDM} \\
\midrule
S01 & — & 43.3 & 29.0 & 27.7 & 29.0 & 78.2 & 66.5 \\
S02 & — & 42.0 & 40.0 & 30.3 & 34.7 & 71.9 & 63.7 \\
S03 & — & 36.7 & 38.3 & 31.3 & 38.7 & 49.5 & 61.1 \\
S04 & — & 43.3 & 37.3 & 41.7 & 46.7 & 68.1 & 55.1 \\
S05 & — & 35.7 & 33.7 & 39.7 & 41.7 & 63.6 & 46.3 \\
S06 & — & 35.7 & 32.3 & 43.0 & 34.3 & 63.6 & 36.3 \\
S07 & — & 35.7 & 31.7 & 33.0 & 31.0 & 83.7 & 63.3 \\
S08 & — & 37.0 & 35.0 & 33.3 & 30.0 & 57.1 & 54.0 \\
S09 & — & 29.7 & 23.3 & 32.0 & 24.7 & 50.8 & 39.6 \\
S10 & — & 29.3 & 33.0 & 22.3 & 40.0 & 63.8 & 44.6 \\
S11 & — & 45.0 & 41.0 & 36.7 & 43.7 & 57.4 & 43.2 \\
S12 & — & 48.7 & 41.7 & 45.7 & 43.3 & 58.6 & 56.3 \\
S13 & — & 29.3 & 37.7 & 33.7 & 33.3 & 50.8 & 28.5 \\
S14 & — & 31.7 & 45.3 & 30.7 & 34.3 & 62.1 & 53.8 \\
S15 & — & 28.0 & 40.7 & 26.7 & 33.3 & 84.2 & 59.1 \\
S16 & — & 41.7 & 33.0 & 40.7 & 43.0 & 49.0 & 33.5 \\
S17 & — & 30.0 & 27.0 & 21.3 & 28.7 & 42.9 & 34.7 \\
S18 & — & 36.0 & 27.0 & 36.7 & 34.0 & 58.1 & 34.9 \\
S19 & — & 33.3 & 31.7 & 28.3 & 38.0 & 48.1 & 35.0 \\
S20 & — & 38.0 & 35.7 & 32.3 & 31.0 & 41.3 & 40.1 \\
\midrule
\textbf{Mean} & \textbf{—} & \textbf{36.5} & \textbf{34.7} & \textbf{33.3} & \textbf{35.7} & \textbf{60.1} & \textbf{47.5} \\
\textbf{95\,\% CI} & — & [34.1,39.2] & [32.4,37.2] & [30.5,36.5] & [33.2,38.2] & [54.7,65.6] & [42.5,52.8] \\
\bottomrule
\end{tabular}
\vspace{0.3em}
{\footnotesize EEGNet = EEGNet (Lawhern 2018)~\cite{lawhern2018eegnet}; Enc = encoder\_ft\_cls; AC-SSM = action-conditioned SSM (this work); Enc+WM = MLP world-model planner; Abl = ablation; CSP = CSP+LDA; MDM = Riemannian MDM.}
\end{table}

\subsection{Cross-Dataset Validation: BCI Competition IV Dataset 2a}
\label{sec:bci2a}

To test whether the discriminability gap generalises beyond PhysioNetMI and to quantify the cost of domain mismatch, we evaluate four conditions on BCI Competition IV Dataset 2a~\cite{brunner2008bci} (BNCI2014\_001): 9 subjects, 4-class MI (left hand, right hand, feet, tongue), 576 trials per subject, 250\,Hz.
Conditions: CSP+LDA, Riemannian MDM, EEGNet~\cite{lawhern2018eegnet}, and the JEPA encoder pretrained on PhysioNetMI (cross-domain transfer) — all evaluated under the same leakage-safe 5-fold LOSO protocol.

Table~\ref{tab:bci2a} reports the results.
Four findings emerge.
First, the CSP--EEGNet gap on BCI2a ($58.7\%$ vs.\ $36.0\%$, $\Delta=22.7$\,pp) closely mirrors the gap on PhysioNetMI ($60.1\%$ vs.\ $38.4\%$, $\Delta=21.7$\,pp), confirming the gap is structural and data-limited.
Second, the JEPA encoder pretrained on PhysioNetMI achieves only $27.9\%$ on BCI2a — statistically indistinguishable from random initialisation ($27.5\%$, $p=0.43$) and from chance ($25.0\%$, $p=0.004$, barely significant at $+2.9$\,pp).
Third, EEGNet outperforms the cross-domain JEPA by $8.1$\,pp ($p=0.004$): supervised training from scratch on the target domain beats a self-supervised encoder trained on a different domain.
Fourth, both methods remain far below CSP ($>20$\,pp gap), confirming the data-limited bottleneck is the binding constraint regardless of training objective or domain.

\begin{table}[t]
\centering
\caption{BCI2a LOSO Balanced Accuracy (9 Subjects, 5 Folds, 4-Class MI, Chance 25.0\,\%)}
\label{tab:bci2a}
\begin{tabular}{lrrr}
\toprule
\textbf{Condition} & \textbf{Mean (\%)} & \multicolumn{2}{c}{\textbf{95\,\% CI (\%)}} \\
 & & Lo & Hi \\
\midrule
CSP + LDA            & 58.7 & 49.1 & 68.7 \\
MDM (Riemannian)     & 50.4 & 40.3 & 60.6 \\
EEGNet~\cite{lawhern2018eegnet} & 36.0 & 31.5 & 41.7 \\
JEPA encoder (cross-domain\textsuperscript{†}) & 27.9 & 26.4 & 29.4 \\
Random encoder (ablation) & 27.5 & 25.9 & 29.1 \\
\midrule
Chance               & 25.0 & — & — \\
\bottomrule
\end{tabular}
\vspace{0.2em}
{\footnotesize \textsuperscript{†}Pretrained on PhysioNetMI, fine-tuned on BCI2a (cross-domain transfer). JEPA vs.\ ablation $p=0.43$ (Wilcoxon); EEGNet vs.\ JEPA $p=0.004$.}
\end{table}

\begin{table}[t]
\centering
\caption{Per-Subject BCI2a Balanced Accuracy (\%)}
\label{tab:bci2a_persubject}
\setlength{\tabcolsep}{4pt}
\begin{tabular}{crrrr}
\toprule
\textbf{Subj.} & \textbf{CSP} & \textbf{MDM} & \textbf{EEGNet} & \textbf{JEPA\textsuperscript{†}} \\
\midrule
S01 & 61.6 & 60.4 & 42.7 & 30.3 \\
S02 & 57.1 & 35.6 & 33.0 & 29.0 \\
S03 & 79.2 & 67.7 & 52.5 & 26.7 \\
S04 & 40.5 & 40.3 & 31.7 & 27.4 \\
S05 & 38.4 & 26.9 & 29.0 & 24.0 \\
S06 & 41.7 & 35.6 & 28.0 & 26.6 \\
S07 & 62.1 & 53.8 & 39.4 & 32.1 \\
S08 & 79.0 & 72.4 & 38.4 & 28.3 \\
S09 & 69.1 & 60.9 & 29.0 & 26.5 \\
\midrule
\textbf{Mean} & \textbf{58.7} & \textbf{50.4} & \textbf{36.0} & \textbf{27.9} \\
\textbf{95\,\% CI} & [49.1,68.7] & [40.3,60.6] & [31.5,41.7] & [26.4,29.4] \\
\bottomrule
\end{tabular}
\vspace{0.2em}
{\footnotesize \textsuperscript{†}JEPA encoder pretrained on PhysioNetMI, fine-tuned on BCI2a (cross-domain).}
\end{table}

\subsection{Decoder Uncertainty Output}
The IntentDecoder produces per-DOF uncertainty $\bm\sigma$ alongside each motor command.
Evaluated over all 1,277 trials from 20 subjects using the final checkpoint, $\sigma_\text{mean} = 0.577 \pm 0.005$ (mean $\pm$ SD across trials).
The near-constant standard deviation ($\pm 0.005$) indicates the decoder has converged to a near-fixed uncertainty estimate regardless of trial content — a failure of calibration.
A weak but statistically significant Pearson correlation exists between per-trial $\sigma$ and prediction error ($r = 0.088$, $p = 0.002$, $N=1277$), confirming a marginal directional signal, but with $r^2 < 0.01$ the uncertainty estimate explains under 1\,\% of error variance.
This finding motivates replacing the fixed-uncertainty decoder with one that adapts to signal quality (e.g.\ via dropout-based uncertainty or an explicit noise-model head).
Despite its current miscalibration, $\bm\sigma$ remains structurally connected to the pipeline: it sets the Kalman filter measurement noise ($R = \text{diag}(\bm\sigma^2)$) and the CEM exploration radius, so improving calibration would directly improve both filter and planner behaviour.

\subsection{Static Classification Does Not Distinguish World-Model Variants}

The MLP world-model condition (encoder\_ft\_cls\_planned) and the AC-SSM both achieve static LOSO accuracy comparable to the JEPA encoder baseline — and all three are $\approx\!22$\,pp below CSP+LDA.
This is expected: static single-trial classification treats each 1\,s window in isolation, providing no temporal motor context.
The MLP transition model introduces a small perturbation that can blur class boundaries (accounting for its marginal accuracy cost vs.\ the base encoder); the AC-SSM's reconstruction and alignment losses regularise against this.

The closed-loop simulation (Section~\ref{sec:closed_loop}) is the correct evaluator for world-model variants: there, the AC-SSM achieves $33.8\%$ convergence versus $0\%$ for the MLP transition model ($p{<}0.0001$), demonstrating that moving action conditioning into the SSM recurrence — rather than bolting on an MLP — is the critical architectural change.

\subsection{Closed-Loop Virtual Arm Simulation}
\label{sec:closed_loop}

We evaluate the world model in its intended setting via a four-controller virtual 6-DOF arm simulation designed to separate decoder quality from architecture quality.
A virtual arm with Euler dynamics $\mathbf{p}_{t+1} = \text{clip}(\mathbf{p}_t + \Delta t \cdot \mathbf{c}_t, -1, 1)$ and $\Delta t = 1/256$\,s is driven for up to $S=384$ steps toward a class-specific DOF target (e.g.\ left-hand $\to [1,0,0,0,0,0]^\top$); convergence is declared when $\|\mathbf{p}_t - \mathbf{p}_\text{goal}\|_2 < 0.25$.

Four controllers are compared on 133 trials (10 subjects, 20\,\% held-out test fold):
\begin{itemize}
  \item \textbf{Oracle}: $\mathbf{c}_t = \mathbf{p}_\text{goal}$ — the exact target DOF issued every step. Represents a perfect decoder (100\,\% accuracy). Establishes whether the arm \emph{can} physically converge.
  \item \textbf{CSP}: $\mathbf{c}_t = \mathbf{p}_\text{goal}^{\hat{y}}$ where $\hat{y}$ is the CSP+LDA prediction on the test window. Represents a classical baseline decoder ($\approx\!60\%$ accuracy in LOSO evaluation).
  \item \textbf{JEPA}: $\mathbf{c}_t = \text{dec}(T(\mathbf{h}_t, \text{dec}(\mathbf{h}_t).\bm\mu)).\bm\mu$ — the MLP-transition world-model decoder ($\approx\!36\%$ static accuracy; near-zero DOF magnitude due to unsupervised latent space).
  \item \textbf{AC-SSM}: at each step, re-encodes the EEG window with the previous command $\mathbf{a}_{t-1}$ as action context, classifies via the linear head fine-tuned on the training fold, and looks up $\mathbf{p}_\text{goal}^{\hat{y}}$ — the same intent-vector lookup used by CSP. This removes dependence on continuous-DOF decoder magnitude while demonstrating that action-conditioned representations produce accurate class predictions.
\end{itemize}

\textbf{Finding (Table~\ref{tab:closedloop}).}
The oracle establishes physical feasibility at $100\%$ convergence and TTT\,=\,193 steps (theoretical minimum: $0.75/\Delta t = 192$ steps).
The CSP controller converges on $26.3\%$ of trials; single-fold per-subject test sets carry high variance (14 trials/subject), pulling the rate below the LOSO 5-fold mean.
On converging CSP trials, TTT\,$\approx$\,193 steps — identical to the oracle — because CSP issues the exact target direction when correct.
The JEPA decoder converges on 0/133 trials ($p{<}0.0001$ vs.\ CSP, Wilcoxon), consistent with near-zero DOF command magnitude throughout all 384 steps.

\begin{table}[t]
\centering
\caption{Four-Controller Closed-Loop Simulation (133 trials, 10 subjects, $S{=}384$, $\varepsilon{=}0.25$)}
\label{tab:closedloop}
\begin{tabular}{lrrrr}
\toprule
\textbf{Metric} & \textbf{Oracle} & \textbf{CSP} & \textbf{JEPA} & \textbf{AC-SSM} \\
\midrule
Convergence rate      & 100.0\,\% & 26.3\,\% & 0.0\,\% & 33.8\,\% \\
Time-to-target (mean) & 193 & 334 & 384 & 319 \\
Final error (L2, DOF) & 0.246 & 1.252 & 0.984 & 1.147 \\
\midrule
\multicolumn{4}{l}{\textit{Wilcoxon (CSP vs.\ JEPA on TTT)}} \\
\quad CSP vs.\ JEPA ($p$) & \multicolumn{4}{r}{$<0.0001$} \\
\quad AC-SSM vs.\ JEPA ($p$) & \multicolumn{4}{r}{$<0.0001$} \\
\bottomrule
\end{tabular}
\vspace{0.2em}
{\footnotesize Oracle TTT=193 matches theoretical minimum (192 steps). CSP and AC-SSM convergence on single fold (14 trials/subject); LOSO 5-fold averages would yield higher rates. AC-SSM uses classifier+intent-lookup identical to CSP; convergence gap vs. oracle reflects static classification accuracy under the data-limited bottleneck.}
\end{table}

Three findings emerge.
First, the oracle establishes physical feasibility: $100\%$ convergence at TTT\,$=\,193$ steps, matching the theoretical minimum of $0.75\,/\,\Delta t = 192$ steps for this target geometry.
Second, the standard JEPA decoder converges on $0/133$ trials — near-zero DOF command magnitude persists throughout all 384 steps, producing $p{<}0.0001$ vs.\ CSP.
Third, and centrally, the AC-SSM classifier-decoder achieves $33.8\%$ convergence (TTT\,=\,319\,$\pm$\,90 steps), significantly above the JEPA baseline ($p{<}0.0001$, Wilcoxon) and statistically indistinguishable from CSP ($26.3\%$, $p=0.21$).
The AC-SSM's TTT is $14.3$ steps faster than CSP on converging trials, consistent with action-conditioned representations focusing earlier on the correct class direction.

These results demonstrate that injecting motor context into the SSM recurrence converts a non-functional decoder into one competitive with the classical spatial-filter baseline, despite operating on the same data-limited 55 labelled trials per subject.
The convergence rate of $33.8\%$ reflects the AC-SSM's static classification accuracy in the closed-loop setting; the gap to the oracle ($100\%$) is the data-limited decoder bottleneck, not an architectural limit.

%% -----------------------------------------------------------------------
\section{Discussion}
\label{sec:discussion}

\subsection{Why Static Classification Understates the World Model's Contribution}
\label{sec:why_fail}

The JEPA objective minimises $\|\text{pool}(f_\theta(\mathbf{X}_\text{ctx})) - f_{\bar\theta}(\mathbf{X}_\text{tgt})\|_2^2$.
This incentivises the encoder to capture temporal autocorrelation — structure that is predictive of future EEG but not necessarily class-discriminative.
Motor imagery discriminability requires ERD signatures to occupy separable regions of the latent space, which is a distinct geometric requirement.

The dominant variance in resting EEG (alpha rhythm fluctuations, electrode drift, movement artefacts) is not motor-class-correlated~\cite{blankertz2011single}.
A temporal JEPA encoder embeds these nuisance factors in its principal components; a linear probe cannot separate the $\ll 1$\,pp class-specific signal from dominant noise.
CSP explicitly resolves this by constructing a basis that maximises the class-to-total variance ratio, which is why the $\approx\!22$\,pp gap ($60.1\%$ vs.\ $\{38.4, 36.6\}\%$ on PhysioNetMI; $58.7\%$ vs.\ $36.0\%$ on BCI2a) is structurally consistent across datasets and training objectives.

Crucially, the AC-SSM does not close this static classification gap — and is not designed to.
The AC-SSM's action conditioning operates in the closed-loop temporal domain, where the previous motor command provides context about which neural trajectory the user is currently on.
In a single-trial LOSO evaluation (each window treated independently), this context is absent.
The AC-SSM's LOSO accuracy is thus comparable to the JEPA baseline, as expected.
The closed-loop simulation (Section~\ref{sec:closed_loop}) is the correct evaluation: there, the AC-SSM achieves $33.8\%$ convergence versus $0\%$ for the JEPA decoder — a qualitative change, not a marginal one.

\subsection{Why the AC-SSM Succeeds Where the MLP Transition Model Failed}
\label{sec:world_model_value}

The original MLP transition model $T(\mathbf{h}, \mathbf{a}) \to \mathbf{h}'$ was trained with a self-consistency loss on a latent space shaped solely by EEG temporal autocorrelation.
That space has no motor semantics: the MLP learns to be self-consistent with near-zero motor commands because the JEPA decoder, operating on an unsupervised latent space, maps to near-zero 6-DOF outputs.
The planner optimises action sequences through a world model that is self-consistent with noise, producing $0\%$ closed-loop convergence.

The AC-SSM resolves this by moving action conditioning upstream into the SSM recurrence itself.
The SiLU gate learns to amplify the motor command injection when the EEG window is consistent with intentional movement and suppress it during rest — a content-adaptive filter grounded in the SSM's own representation.
The EEG reconstruction loss provides a label-free training signal ($\approx 56{,}000$ timestep-level gradients vs.\ 55 classification labels per fold), shaping the encoder toward brain dynamics rather than static trial features.
The result is a latent space where motor class directions are structurally encoded — enabling the closed-loop classifier-decoder to issue correct DOF commands on $33.8\%$ of trials from 55 training examples.

The world model's three functions remain intact:
\begin{enumerate}
  \item \textbf{Closed-loop prediction}: the action-conditioned SSM recurrence predicts how the latent brain state evolves after a motor command, enabling the planner to roll out trajectories without waiting for a new EEG observation.
  \item \textbf{Latency minimisation}: the CEM planner (Algorithm~\ref{alg:cem}) with $\gamma=0.8$ values convergence at step~1 by $2.8\times$ over step~3, directly reducing command latency.
  \item \textbf{Uncertainty-proportional exploration}: CEM samples from $\mathcal{N}(\bm\mu_\text{dec}, \bm\sigma_\text{dec})$; high uncertainty drives broader search before commitment.
\end{enumerate}

No classical comparator (CSP, MDM) provides any of these functions.
The four-controller simulation confirms the architecture's viability: oracle $100\%$, AC-SSM $33.8\%$, CSP $26.3\%$ ($p=0.21$, equivalent), JEPA $0\%$ ($p{<}0.0001$).

\subsection{Why Cross-Modal CNS Pretraining Did Not Help}
\label{sec:cns}

The CNS encoder was trained for only 3 epochs on synthetic MC\_Maze data (real NLB data unavailable without DANDI credentials), producing a frozen teacher that encodes spike-train autocorrelation rather than arm-kinematic motor manifold geometry.
A properly trained CNS encoder (Safaie et al.~\cite{safaie2023preserved} report motor manifold preserved across species) should provide a motor-class-discriminative teacher.
The negative result here reflects the quality of the CNS teacher, not the architectural principle.

\subsection{Predictability vs.\ Discriminability: General Principle}
\label{sec:pred_disc}

Self-supervised temporal prediction learns the first-order statistics of the data manifold — structure useful for compression and generation, not classification.
Class-discriminative structure requires class-specific features to lie in the \emph{high-variance} directions of the representation.
For EEG-MI, this is guaranteed only under a discriminative objective (CSP, contrastive learning with class labels).
This principle extends beyond EEG: any domain where the dominant signal variance is not class-correlated will exhibit the same pattern.
The contribution of this work is to establish the principle quantitatively across two independent datasets: both JEPA (predictive) and EEGNet (discriminative) fall $\approx\!22$\,pp below CSP, and neither distinguishably outperforms random initialisation at $\leq\!576$ trials per subject.

\subsection{Limitations and Future Work}
\label{sec:limitations}

\begin{enumerate}
  \item \textbf{No within-domain BCI2a JEPA pretraining.} A BCI2a-native JEPA encoder (pretrained on BCI2a EEG, without labels) was computationally infeasible to evaluate at full scale (30 epochs on 83,022 windows $\approx$54\,h CPU). The cross-domain transfer result ($27.9\% \approx$ ablation) establishes an important lower bound; a within-domain BCI2a JEPA is expected to close the $8.1$\,pp gap vs.\ EEGNet and remains future work.
  \item \textbf{No EEGConformer or BENDR baseline.} EEGNet~\cite{lawhern2018eegnet} is included and matches the JEPA encoder. EEGConformer~\cite{song2022eegconformer} and fine-tuned BENDR~\cite{kostas2022bendr} are not in the current LOSO harness; given that even the discriminative EEGNet cannot close the gap, these are unlikely to change the conclusion at this trial count.
  \item \textbf{Closed-loop simulation uses a static EEG window.} The virtual arm simulation re-encodes the same 1\,s EEG window at each step with the updated $\mathbf{a}_\text{prev}$, rather than streaming new EEG.
    In real deployment, the encoder would receive a genuinely updated window as the user's neural state evolves — expected to improve convergence further.
    A hardware-in-the-loop evaluation with an online EEG stream is the next step.
  \item \textbf{No continuous-DOF decoder.} The AC-SSM closed-loop controller uses a discrete classifier $\to$ intent lookup, identical to CSP.
    A continuous 6-DOF decoder (replacing the lookup with a regression head trained on intent vectors) would enable smooth trajectories; this requires either more labelled data or a stronger discriminative pretraining objective to resolve the decoder magnitude problem.
  \item \textbf{Dataset size.} 69 trials/subject is below the threshold at which fine-tuned neural encoders reliably outperform CSP~\cite{kostas2022bendr}, explaining the 22.7\,pp gap.
  \item \textbf{CPU inference speed.} \texttt{decode\_step()} runs at $26.9$\,ms ($37$\,Hz) and \texttt{self\_condition()} at $10.9$\,ms on CPU. This supports typical BCI control loop rates ($10$--$64$\,Hz) but leaves no headroom for higher rates. CUDA compilation hangs on the sequential ZOH-SSM scan with RTX\,5070 (compute capability 12.x); diagonal S4D parameterisation would resolve this and is expected to reduce latency by one to two orders of magnitude.
\end{enumerate}


%% -----------------------------------------------------------------------
\section{Conclusion}
\label{sec:conclusion}

We presented the Action-Conditioned ZOH-SSM World Model, a BCI architecture that injects the previous motor command into the SSM recurrence via a SiLU-gated adaptive bias, trained jointly on cross-entropy classification, unsupervised EEG reconstruction, and cosine alignment to a JEPA baseline.
A hybrid CEM+gradient MPC planner exploits the world model for latency-minimising closed-loop control.

On static classification, the AC-SSM matches the JEPA encoder baseline under the same data-limited conditions (20 subjects, $\approx\!69$ trials/subject), with both architectures $\approx\!22$\,pp below CSP+LDA across two independent datasets — establishing that the bottleneck is labelled data volume, not the training objective.

In closed-loop virtual arm simulation, the AC-SSM achieves $33.8\%$ target convergence versus $0\%$ for the standard JEPA decoder ($p{<}0.0001$, Wilcoxon) and $26.3\%$ for CSP ($p=0.21$, not significant).
This is the first demonstration that action conditioning built into an SSM recurrence — rather than added as a bolted-on MLP transition model — can produce functional closed-loop prosthetic control from 55 labelled EEG trials per subject, with performance competitive with classical spatial-filter decoders.

The architecture, leakage-free evaluation protocol, and full result data are released to provide a reproducible foundation for action-conditioned world-model research in neural BCI.


%% -----------------------------------------------------------------------
\section*{Ethical Statement}
This study uses only publicly available, fully de-identified data (PhysioNet EEG Motor Movement/Imagery Dataset, CC-BY 4.0).
No human subjects were recruited; institutional ethics review was not required.

\section*{Acknowledgment}
The author thanks the PhysioNet team for open data stewardship and the developers of MNE-Python, pyriemann, and PyTorch.
No funding was received for this work.

\section*{Data Availability}
The PhysioNet Motor Imagery dataset used in this study is freely available under the Creative Commons Attribution 4.0 International licence at \url{https://physionet.org/content/eegmmidb/}.
The full evaluation codebase, pre-trained checkpoints, and per-subject result files are available at \url{https://github.com/josephlw4/noosphere} (to be made public upon acceptance).

%% -----------------------------------------------------------------------
\appendix
\section{Implementation Details}
\label{app:impl}

\subsection{ZOH-SSM Hyperparameters}
\begin{itemize}
  \item Layers $L=4$; state dim $d_\text{state}=64$; model dim $d_\text{model}=128$; dropout $p=0.1$.
  \item \texttt{torch.compile(dynamic=True)} at module level; cold-compile $\sim$254\,s, subsequent folds $\sim$3.5\,s (cache hit).
\end{itemize}

\subsection{JEPA Pretraining}
\begin{itemize}
  \item Context fraction 0.7, target 0.3; EMA $m=0.996$; batch 32; Adam $\text{lr}=3\times10^{-4}$; 50 epochs.
\end{itemize}

\subsection{Transition Model}
\begin{itemize}
  \item $d_\text{hidden}=256$; GELU + LayerNorm; skip $\mathbf{W}_\text{skip}$ initialised to identity; self-consistency weight 0.1.
  \item Trained jointly with supervised loss (10 fine-tuning epochs, AdamW $\text{lr}=10^{-3}$).
\end{itemize}

\subsection{LatencyPlanner}
\begin{itemize}
  \item $H=3$; $\gamma=0.8$; $K=32$ MC rollouts; elite fraction $\rho=0.25$; $n_\text{iter}=12$ gradient steps; Adam $\text{lr}=0.05$.
  \item Fast path: single forward pass $T(\mathbf{h}, \text{dec}(\mathbf{h}).\bm\mu)$, $10.9$\,ms CPU ($N=2560$ trials, post-compile). \texttt{decode\_step()}: $26.9$\,ms, $37$\,Hz.
  \item \texttt{torch.enable\_grad()} inside \texttt{plan()} to work within \texttt{torch.no\_grad()} inference contexts.
\end{itemize}

\subsection{LOSO Command (v4 — final)}
\begin{verbatim}
.venv/bin/python -m \
  v2_digital_self_replication.cli.eval_loso \
  --subjects 1-20 \
  --classes left_hand right_hand feet \
  --folds 5 --ft-epochs-cls 5 \
  --fast --eegnet --eegnet-epochs 50 \
  --device cpu \
  --checkpoint .../supervised_best.pt \
  --cns-checkpoint .../jepa_encoder_cns_pretrained.pt \
  --output .../loso_results_v4.json
\end{verbatim}


%% -----------------------------------------------------------------------
\bibliographystyle{IEEEtran}
\bibliography{references}

\end{document}
